\chapter*{TD2 : Convergence de variables aléatoires et théorèmes limites.}

\setcounter{exercice}{0}

\begin{mdframed}
\begin{exercice}
On reprend l'exercice 1 de la feuille précédente. On rappelle la fonction de masse de $X$ :

\begin{align*}
\mathbb{P}(X = -2) = \frac{1}{30}, \\
\mathbb{P}(X = 0) = \frac{1}{20}, \\
\mathbb{P}(X = 1) = \frac{1}{6}, \\ 
\mathbb{P}(X = 3) = \frac{1}{12}, \\ 
\mathbb{P}(X = 5) = \frac{2}{3}.
\end{align*}

\begin{enumerate}
\item Montrer que $Y = -2X + 1$ est une v.a. discrète et déterminer sa fonction de masse. 
\item Même question pour $Z = Y^2$.
\end{enumerate}
\end{exercice}
\end{mdframed}

\medskip

\begin{mdframed}
\begin{exercice}
On reprend l'exercice 4 de la feuille précédente. On rappelle la fonction de répartition de $X$ :

$$
F_X(x) = \begin{cases} 
0 & \text{si } x < 0, \\ 
\frac{x}{4} & \text{si } 0 \leqslant x < 1, \\ 
\frac{1}{4} & \text{si } 1 \leqslant x < 3, \\ 
\frac{3x-7}{8} & \text{si } 3 \leqslant x < 5, \\ 
1 & \text{si } x \geqslant 5. 
\end{cases}
$$

\begin{enumerate}
\item Justifier que la variable aléatoire $Y = 1/X$ est bien définie. 
\item Montrer que $Y$ est à densité et calculer sa densité.
\end{enumerate}
\end{exercice}
\end{mdframed}

\medskip

\begin{mdframed}
\begin{exercice}
Soit $f_X$ la densité d'une variable aléatoire réelle $X$ telle que $\mathbb{P}(X \geq 0) > 0$. On pose 

$$
f_Y(t) = af_X(|t|), \quad t \in \mathbb{R}
$$

\begin{enumerate}
\item Calculer $a$ (en fonction de $f_X$) pour que $f_Y$ soit la densité d'une variable aléatoire $Y$. 
\item Montrer que $Y$ et $-Y$ ont même loi.
\end{enumerate}
\end{exercice}
\end{mdframed}

\medskip

\begin{mdframed}
\begin{exercice}
Soit $X$ une variable aléatoire de densité $f$ continue et de fonction de répartition $F$. 

\begin{enumerate}
\item Soit $Z = 2X\mathbb{1}_{\{X<0\}} + 3X\mathbb{1}_{\{X \geq 0\}}$. Donner la fonction de répartition de $Z$ en fonction de $F$, puis sa densité en fonction de $f$. 
\item Soit $Y = X^2$. Donner la fonction de répartition de $Y$ en fonction de $F$, puis sa densité en fonction de $f$.
\end{enumerate}
\end{exercice}
\end{mdframed}

\medskip

\begin{mdframed}
\begin{exercice}
Soit $X$ une variable aléatoire de fonction de répartition $F$, et soit $A = ]a, b]$ un intervalle. 

\begin{enumerate}
\item Soit $Y = \mathbb{1}_A(X)$. Donner la fonction de répartition de $Y$ en fonction de $F$. (On pourra observer que $Y \sim \text{Ber}(p)$ pour une valeur de $p$ à déterminer). 
\item On suppose de plus $a > 0$. Soit $Z = X\mathbb{1}_A(X)$. Donner la fonction de répartition de $Z$ en fonction de $F$.
\end{enumerate}
\end{exercice}
\end{mdframed}

\medskip

\begin{mdframed}
\begin{exercice}
Soit $X$ une variable aléatoire de fonction de répartition $F$. Soit $(a, b) \in \mathbb{R}^2$. 

\begin{enumerate}
\item Exprimer à l'aide de $F$ la fonction de répartition de la variable aléatoire $aX + b$. 
\item Sous quelle condition sur le couple $(a, b)$ la fonction $t \mapsto F(at + b)$ est-elle une fonction de répartition ? Sous cette condition, de quelle variable aléatoire cette fonction est-elle la fonction de répartition ?
\end{enumerate}
\end{exercice}
\end{mdframed}

